% !TEX program = xelatex
% Research draft based on sadekovlar/misis-article-latex.
% The original template is identified in TEMPLATE.md.
\documentclass{misisarticle}

% Уточнить классификационные коды и аффилиацию перед сдачей.
% Пустые коды не выводятся; коды из демонстрационной статьи не перенесены.
\grnti{}
\asjc{}
\oecd{}

% Рабочая формулировка: согласовать с научным руководителем.
\title{Внешняя калибровка нескольких ToF-камер с использованием калибровочного объекта}

\misisauthor{Бауэр Вольдемар Владимирович}{Группа БПМ-23-ПО-1}{}{}
% Руководитель обозначен как руководитель, а не автоматически как соавтор.
\misisauthor{Мишуров С.~С.}{Научный руководитель}{}{}


\setabstract{Рассматривается задача определения взаимного положения нескольких камер времени пролёта по наблюдениям калибровочного объекта известной геометрии. Обсуждаются методы калибровки камер, оценки позы по соответствиям трёхмерных и двумерных точек и регистрации облаков точек. Формулируется подход, в котором начальные оценки внешних параметров получают по изображениям многогранной мишени, а затем уточняют совместной оптимизацией. При наличии согласованных измерений глубины предлагается дополнить функцию ошибки расстояниями до граней мишени. Определены требования к данным, условия наблюдаемости и критерии последующей проверки. Работа представляет постановку исследования и описание предполагаемого метода; экспериментальные результаты на данном этапе не приводятся.}

\setkeywords{ToF-камера, внешняя калибровка, калибровочный объект, PnP, совместная оптимизация, карта глубины}

\hypersetup{
  pdftitle={Внешняя калибровка нескольких ToF-камер с использованием калибровочного объекта},
  pdfauthor={Бауэр Вольдемар Владимирович},
  pdfsubject={НИР: постановка задачи и предполагаемый метод}
}

\begin{document}
\maketitle
\section{Введение}

Камеры времени пролёта (time-of-flight, ToF) позволяют получать измерения расстояния до наблюдаемой поверхности. Несколько таких камер могут использоваться для построения общего трёхмерного представления сцены. Однако карты глубины и облака точек первоначально заданы в отдельных системах координат. Для их объединения необходимо определить вращения и переносы между камерами, то есть выполнить внешнюю калибровку.

Эта задача отличается от внутренней калибровки, в которой оценивают параметры проекции и искажения изображения. Даже при известных внутренних параметрах неверная взаимная ориентация камер приводит к несовпадению наблюдений одной поверхности. При этом ошибки ToF-измерений могут зависеть от расстояния, свойств поверхности и условий получения данных~\cite{foix2011}. Поэтому оценку внешних параметров следует отделять от оценки качества самой глубины.

В данной работе рассматривается калибровка по объекту известной геометрии, например по кубу с различимыми элементами на гранях. Такой объект задаёт общую систему координат для наблюдений нескольких камер. Камеры могут видеть разные элементы мишени, если их координаты определены в одной системе и взаимное положение граней известно. Это позволяет исследовать конфигурации, для которых поиск соответствий непосредственно между изображениями сцены затруднён.

Цель исследования --- разработать и оценить способ определения взаимного положения ToF-камер по наблюдениям калибровочного объекта. В качестве базового решения рассматривается оценка позы по изображениям с последующим совместным уточнением. Дополнительное направление --- использование глубины и анализ чувствительности результата к погрешностям геометрии мишени. На текущем этапе сформулированы модель, требования к входным данным и план проверки; реализация и экспериментальная оценка остаются последующими этапами.

\section{Анализ существующих работ}

\subsection{Внутренняя калибровка и оценка позы}

Метод Zhang~\cite{zhang2000} использует несколько изображений плоского шаблона для определения внутренних параметров камеры и последующего нелинейного уточнения. Для рассматриваемой задачи такая калибровка является предварительным этапом: параметры проекции должны быть известны либо оценены отдельно. Калибровка отдельной камеры сама по себе не задаёт взаимное положение всех камер установки.

При известных внутренних параметрах положение объекта относительно камеры можно оценивать по соответствиям между его трёхмерными точками и их проекциями. В работе Lepetit и соавторов~\cite{lepetit2009} предложен метод EPnP, вычислительная сложность которого линейна по числу соответствий. Это подходящая основа для получения начальной оценки позы, однако качество результата зависит от распределения наблюдаемых точек и точности их локализации. В прикладной реализации дополнительно необходимы отбраковка ошибочных соответствий и проверка геометрической допустимости решения.

\subsection{Маркеры и многокамерная калибровка}

Система AprilTag~\cite{olson2011} позволяет обнаруживать и идентифицировать искусственные визуальные метки. Идентификаторы помогают сопоставлять точки изображения с заранее заданной геометрией объекта. Однако наличие меток не определяет автоматически координаты их углов в системе всей мишени: расположение меток на гранях должно быть задано и проверено отдельно. Выбор семейства маркеров в настоящей работе будет зависеть от фактической конструкции калибровочного объекта.

Детектирование шаблона в ToF-изображениях требует отдельного внимания. Hansard и соавторы~\cite{hansard2014} рассматривают обнаружение шахматной сетки с использованием её глобальной геометрии. Эта работа показывает значимость специализированной обработки таких изображений; она не даёт оснований заранее считать произвольный детектор маркеров работоспособным на данных ещё не изученной камеры.

Hansard и соавторы~\cite{hansard2015} рассматривают взаимную калибровку нескольких систем ToF и цветных камер с использованием калибровочного шаблона. Эта работа непосредственно относится к объединению разных наблюдений в общую геометрическую модель. Для нашей постановки существенны доступность изображения мишени, известность её геометрии и согласованность систем координат каналов. При этом применение конкретной схемы измерений должно определяться доступными данными установки.

Совместное уточнение параметров камер и структуры наблюдений относится к классу задач bundle adjustment, рассмотренному Triggs и соавторами~\cite{triggs2000}. Здесь структура мишени предполагается известной, а оптимизируемыми величинами служат внешние параметры камер и позы объекта в отдельных наблюдениях. Такая постановка позволяет использовать все наблюдения в одной функции ошибки и явно зафиксировать систему отсчёта.

\subsection{Особенности глубины ToF и регистрация}

Обзор Foix и соавторов~\cite{foix2011} систематизирует принципы работы и калибровки lock-in ToF-камер, а также источники ошибок измерения. Из этого следует необходимость учитывать качество глубины при её включении в геометрическую оптимизацию. В настоящей постановке измерения около границ объекта и недостоверные отсчёты предлагается исключать, а оставшиеся остатки нормировать с учётом принятой модели погрешности.

Другой подход к совмещению измерений --- регистрация трёхмерных данных. Алгоритм ICP Besl и McKay~\cite{besl1992} итеративно уточняет преобразование на основе соответствий между геометрическими элементами. Применительно к нескольким камерам такой подход требует достаточной общей геометрии и подходящего начального приближения. Если камеры видят разные грани мишени, прямое сопоставление облаков может оказаться недостаточно определённым. Поэтому в качестве базового метода здесь выбрана привязка к известной геометрии объекта, а регистрация облаков рассматривается как возможный дополнительный способ проверки или уточнения при наличии перекрытия.

Таким образом, рассматриваемая схема опирается на известные методы, а исследовательский вопрос состоит в их применимости к доступной ToF-установке, устойчивости к неполным наблюдениям и чувствительности к погрешностям мишени. Преимущество использования глубины по сравнению с калибровкой только по изображениям должно быть проверено экспериментально.

\section{Постановка задачи}

\subsection{Входные данные и обозначения}

Пусть установка содержит $N\geq2$ неподвижные камеры, а калибровочный объект наблюдается в $M$ положениях. Под одним положением понимается набор наблюдений, для которых объект и камеры сохраняют взаимное расположение. Для движущейся мишени требуется синхронизация; для неподвижной допустима последовательная съёмка при сохранении её позы.

Для камеры $i$ предполагаются известными матрица внутренних параметров $K_i$ и параметры дисторсии $\boldsymbol\delta_i$. В системе объекта $O$ заданы координаты характерных точек $\mathbf X_j$. Для видимых точек измерены координаты $\mathbf u_{ijs}$ на изображении камеры $i$ в положении объекта $s$. Не требуется, чтобы каждая камера видела все точки или все положения.

Примем систему первой камеры $C_0$ за мировую систему $W$. Обозначим через $T_i\in SE(3)$ преобразование из $W$ в систему камеры $C_i$, а через $P_s\in SE(3)$ --- преобразование из системы объекта в $W$. Тогда $T_0=I$, а модель наблюдения имеет вид
\begin{equation}
  \widehat{\mathbf u}_{ijs}
  =\pi_i\!\left(T_iP_s\widetilde{\mathbf X}_j\right),
  \label{eq:projection}
\end{equation}
где $\widetilde{\mathbf X}_j$ --- однородные координаты точки, а $\pi_i$ включает перспективную проекцию, $K_i$ и модель дисторсии. Для устранения неоднозначности масштаба геометрия объекта должна быть известна в метрических единицах.

\subsection{Искомый результат и условия определимости}

Необходимо оценить $T_i$ для всех камер и вспомогательные позы $P_s$. Преобразование точки из камеры $i$ в систему камеры $j$ после калибровки задаётся матрицей
\begin{equation}
  T_{j\leftarrow i}=T_jT_i^{-1}.
  \label{eq:relative}
\end{equation}
Вместе с матрицами должны сохраняться направление преобразования, единицы длины и сведения об использованных наблюдениях.

Связь камер с наблюдаемыми положениями объекта можно представить двудольным графом. Ребро соответствует наличию допустимых наблюдений элементов мишени. Граф должен быть связным; иначе взаимное положение отдельных компонент не определяется. Связность является необходимым, но не достаточным условием: также нужны информативные расположения точек и позы, позволяющие исключить вырожденные конфигурации и неоднозначность плоской мишени. Для выбранной ниже инициализации дополнительно требуется связная цепочка наблюдений, допускающих самостоятельную оценку позы через PnP.

Сами по себе разные видимые грани не препятствуют калибровке, если их геометрическая связь известна. Однако ошибки размеров объекта, углов между гранями и расположения меток могут систематически смещать оценки. Поэтому номинальная модель куба и фактическая точность его изготовления должны учитываться раздельно.

\subsection{Требования к экспериментальным данным}

До реализации необходимо уточнить число и модели камер; наличие изображений, карт глубины и показателей достоверности; единицы глубины и её физический смысл; внутренние параметры и преобразование между каналами; геометрию и разметку объекта; порядок съёмки и наличие эталонных измерений. В текущем описании эти характеристики не считаются установленными.

Базовая ветвь метода предполагает доступность изображений, на которых различимы опорные элементы объекта. Если имеются только карты глубины или облака точек, потребуется отдельная процедура выделения геометрии мишени, и применение описанного ниже PnP без дополнительных данных будет невозможно. Использование глубины в совместной оптимизации рассматривается как расширение базовой ветви.

\section{Описание предполагаемого метода}

\subsection{Подготовка наблюдений}

Сначала формируется метрическая модель объекта: координаты опорных точек, принадлежность граней и, при наличии маркеров, соответствие их идентификаторов геометрии. Затем на изображениях обнаруживаются видимые элементы, уточняются координаты углов и отбрасываются неоднозначные соответствия. Для каждого набора сохраняются идентификаторы камеры, положения объекта и точки. Пиксельные координаты и модель проекции должны соответствовать одному представлению изображения: нельзя одновременно использовать исправленные изображения и повторно применять к ним исходную дисторсию.

Наблюдения желательно получать при нескольких положениях и ориентациях мишени. Повторение практически одинакового ракурса увеличивает число отсчётов, но не заменяет геометрического разнообразия. Отдельно проверяются достаточность числа точек, их распределение по изображению и видимость объекта перед камерой.

\subsection{Начальная оценка внешних параметров}

Для каждой допустимой пары «камера --- положение объекта» решается задача PnP. Полученная матрица $A_{is}$ переводит координаты объекта в систему камеры, то есть в идеальной модели
\begin{equation}
  A_{is}=T_iP_s.
\end{equation}
EPnP~\cite{lepetit2009} может использоваться для начального решения и поддерживает плоские и пространственные конфигурации точек. При наблюдении только одной грани необходимо отдельно контролировать неоднозначность и устойчивость оценки, проверяя положительность глубин и согласованность с другими наблюдениями. Оценка дополняется устойчивой отбраковкой выбросов и локальным уточнением ошибки проекции.

Если камеры $i$ и $j$ наблюдают одну позу объекта, начальное относительное преобразование вычисляется как
\begin{equation}
  T_{j\leftarrow i}^{(s)}=A_{js}A_{is}^{-1}.
  \label{eq:initial}
\end{equation}
При отсутствии общих наблюдений у пары камер связь может быть восстановлена по пути в связном графе. Полученные оценки служат инициализацией; итоговое решение не должно зависеть только от одного такого пути. Для проверки соглашений о координатах используется обратное проецирование известных точек по формуле~\eqref{eq:projection}.

\subsection{Совместное уточнение по изображениям}

После инициализации внешние параметры и позы мишени уточняются одновременно. Для множества корректных наблюдений $\Omega_I$ определяется функция
\begin{equation}
  E_I=\sum_{(i,j,s)\in\Omega_I}
  \rho_I\!\left(
    \frac{\|\mathbf u_{ijs}-\widehat{\mathbf u}_{ijs}\|^2}
         {\sigma_{u,ijs}^2}
  \right),
  \label{eq:imagecost}
\end{equation}
где $\sigma_{u,ijs}$ характеризует масштаб ошибки локализации, а $\rho_I$ --- робастная функция потерь. При отсутствии индивидуальных оценок точности допускается общий масштаб, который должен быть явно указан. Оптимизация проводится по $T_i$ и $P_s$ с фиксированным $T_0=I$; вращения параметризуются с сохранением принадлежности $SO(3)$. Внутренние параметры и геометрия объекта в базовой постановке остаются фиксированными.

Такая задача является специализированным вариантом совместной оптимизации~\cite{triggs2000}. Робастная функция снижает влияние отдельных неверных наблюдений, но не устраняет систематическую ошибку геометрии или внутренней калибровки. Поэтому после оптимизации требуется анализ остатков по камерам, положениям и участкам изображения.

\subsection{Дополнение измерениями глубины}

Если доступны надёжные карты глубины, можно использовать точки внутренних областей граней. Для каждого пикселя сначала восстанавливаются координаты в системе датчика глубины. Если значение $z$ означает осевую координату, то для исправленного пикселя $\mathbf v=(u,v,1)^\mathsf T$
\begin{equation}
  \mathbf q=zK_D^{-1}\mathbf v.
\end{equation}
Если же значение $r$ означает расстояние вдоль луча, используется другая формула:
\begin{equation}
  \mathbf q=r\frac{K_D^{-1}\mathbf v}
                       {\|K_D^{-1}\mathbf v\|}.
\end{equation}
Здесь $K_D$ относится к каналу глубины. Перед использованием точку необходимо перевести в выбранную систему $C_i$ по известному преобразованию между каналами. Простое совпадение размеров изображений не гарантирует их геометрического совмещения.

Пусть грань $f$ в системе объекта описывается единичной нормалью $\mathbf n_f$ и уравнением $\mathbf n_f^\mathsf T\mathbf x+b_f=0$. Для измеренной точки $\mathbf q_{iks}$, уже выраженной в $C_i$, остаток равен
\begin{equation}
  r^D_{iks}=\mathbf n_f^\mathsf T
    \operatorname{xyz}\!\left((T_iP_s)^{-1}
      \widetilde{\mathbf q}_{iks}\right)+b_f,
\end{equation}
где $f$ --- установленная принадлежность точки грани, а $\operatorname{xyz}$ извлекает три пространственные координаты. Наборы точек ограничиваются внутренними областями граней, чтобы уменьшить влияние смешанных измерений на границах.

Совместный критерий записывается как
\begin{align}
  E &= E_I+\lambda E_D,\\
  E_D &=\sum_{(i,k,s)\in\Omega_D}
    \rho_D\!\left(\frac{(r^D_{iks})^2}{\sigma_{d,iks}^2}\right),
\end{align}
где $\sigma_{d,iks}$ --- масштаб погрешности расстояния до плоскости, а $\lambda$ управляет вкладом глубины. Этот масштаб должен учитывать распространение погрешности измерения вдоль луча на нормаль грани. Нормирование делает остатки безразмерными и позволяет согласовать их вклад. Влияние числа и пространственной корреляции точек также следует учитывать при выборе выборки и веса. Значение $\lambda=0$ соответствует базовой калибровке по изображениям.

Измерения одной плоскости не ограничивают одинаково все компоненты позы. Поэтому глубина дополняет визуальные соответствия, а достаточность геометрии оценивается для набора наблюдений в целом. Конкретная модель достоверности и настройки робастной функции должны выбираться после изучения характеристик камеры, а не назначаться как универсальные константы.

\subsection{Критерии последующей проверки}

Сопоставление базового и расширенного вариантов предполагается проводить на одинаковых наборах наблюдений. Следует анализировать ошибку репроекции на отложенных положениях мишени, устойчивость оценок при изменении состава кадров и согласованность независимо оценённых парных преобразований по циклам графа до совместной оптимизации. Последняя является внутренней проверкой и сама по себе не доказывает абсолютную точность. Циклы, составленные только из итоговых матриц $T_jT_i^{-1}$, замыкаются тождественно и не являются критерием качества.

При проверке на отложенных положениях внешние параметры камер фиксируются. Для каждого положения оценивается одна общая поза объекта по части наблюдений, а ошибка вычисляется на оставшейся части. Независимая подгонка позы для каждой камеры не проверяла бы их взаимную калибровку.

При наличии независимого эталона можно измерять ошибки вращения и переноса. Без такого эталона допустимо сообщать ошибки проекции, геометрические остатки и повторяемость, но нельзя отождествлять их с истинной погрешностью внешних параметров. Проверка на глубине должна по возможности использовать данные, не участвовавшие в оптимизации.

Отдельный предполагаемый анализ --- внесение контролируемых возмущений в размеры мишени, углы между гранями и координаты меток. Он позволит оценить чувствительность решения к неточности изготовления и сопоставить её с влиянием ошибок изображения и глубины. Такой анализ пока является планом исследования; утверждений о достигнутой точности и преимуществе конкретной конфигурации в настоящей работе не делается.

\begin{thebibliography}{8}

\bibitem{foix2011}
Foix~S., Alenyà~G., Torras~C. Lock-in Time-of-Flight (ToF) Cameras: A Survey.
\textit{IEEE Sensors Journal}. 2011. Vol.~11, no.~9. P.~1917--1926.
DOI: \href{https://doi.org/10.1109/JSEN.2010.2101060}{10.1109/JSEN.2010.2101060}.

\bibitem{zhang2000}
Zhang~Z. A Flexible New Technique for Camera Calibration.
\textit{IEEE Transactions on Pattern Analysis and Machine Intelligence}.
2000. Vol.~22, no.~11. P.~1330--1334.
DOI: \href{https://doi.org/10.1109/34.888718}{10.1109/34.888718}.

\bibitem{lepetit2009}
Lepetit~V., Moreno-Noguer~F., Fua~P. EPnP: An Accurate $O(n)$ Solution to the PnP Problem.
\textit{International Journal of Computer Vision}.
2009. Vol.~81, no.~2. P.~155--166.
DOI: \href{https://doi.org/10.1007/s11263-008-0152-6}{10.1007/s11263-008-0152-6}.

\bibitem{olson2011}
Olson~E. AprilTag: A robust and flexible visual fiducial system.
\textit{Proceedings of the IEEE International Conference on Robotics and Automation}.
2011. P.~3400--3407.
\url{https://april.eecs.umich.edu/papers/details.php?name=olson2011tags}.

\bibitem{hansard2014}
Hansard~M., Horaud~R., Amat~M., Evangelidis~G.
Automatic Detection of Calibration Grids in Time-of-Flight Images.
\textit{Computer Vision and Image Understanding}.
2014. Vol.~121. P.~108--118.
DOI: \href{https://doi.org/10.1016/j.cviu.2014.01.007}{10.1016/j.cviu.2014.01.007}.

\bibitem{hansard2015}
Hansard~M., Evangelidis~G., Pelorson~Q., Horaud~R.
Cross-calibration of Time-of-flight and Colour Cameras.
\textit{Computer Vision and Image Understanding}.
2015. Vol.~134. P.~105--115.
DOI: \href{https://doi.org/10.1016/j.cviu.2014.09.001}{10.1016/j.cviu.2014.09.001}.

\bibitem{triggs2000}
Triggs~B., McLauchlan~P.~F., Hartley~R.~I., Fitzgibbon~A.~W.
Bundle Adjustment --- A Modern Synthesis.
\textit{Vision Algorithms: Theory and Practice}.
Lecture Notes in Computer Science. Vol.~1883.
Springer, 2000. P.~298--372.
DOI: \href{https://doi.org/10.1007/3-540-44480-7_21}{10.1007/3-540-44480-7\_21}.

\bibitem{besl1992}
Besl~P.~J., McKay~N.~D. A Method for Registration of 3-D Shapes.
\textit{IEEE Transactions on Pattern Analysis and Machine Intelligence}.
1992. Vol.~14, no.~2. P.~239--256.
DOI: \href{https://doi.org/10.1109/34.121791}{10.1109/34.121791}.

\end{thebibliography}

\end{document}
